\documentclass[10pt,a4paper]{amsart}
\usepackage[margin=19mm]{geometry}
\usepackage{amsmath,amssymb,mathtools}
\usepackage[scr=rsfs,cal=euler]{mathalfa}
\usepackage{xfrac}
\usepackage[unicode,hidelinks,bookmarks=false]{hyperref}
\newcommand{\ie}{i.e.\ }
\newcommand{\cf}{cf.\ }
\newcommand{\resp}{respectively\ }
\newcommand{\Subsection}{\subsection}
\providecommand{\nobreakdash}{\nobreak\hbox{-}}
\setlength{\emergencystretch}{2em}
\multlinegap0pt
\usepackage{amssymb}
\usepackage{mathtools}
\usepackage[shortlabels]{enumitem}
\usepackage{xcolor}
\usepackage{tikz}
\newtheorem{lemma}{Lemma}
\title{Hamilton decompositions of the directed 7-torus at odd modulus via root-flat certificates and a prefix-count construction}
\author{SangHyun Park}
\date{Source: arXiv:2605.00660v1 (CC BY 4.0)}
\begin{document}
\maketitle

\noindent\emph{Source and licence.} Hamilton decompositions of the directed 7-torus at odd modulus via root-flat certificates and a prefix-count construction. Version arXiv:2605.00660v1 (CC BY 4.0), distributed by arXiv under the Creative Commons Attribution 4.0 International licence, http://creativecommons.org/licenses/by/4.0/, as recorded in the arXiv licence metadata for this exact version and shown on the abstract page https://arxiv.org/abs/2605.00660v1. Full licence notice: \texttt{licenses/arxiv-2605.00660-CC-BY-4.0.txt}.\par\smallskip

\noindent\textit{Editorial scope.} Counted unit: the article's lemma lem:count-criterion (source0.tex lines 406--425). The article's complete original proof of it is delivered with it (source0.tex lines 427--478, one \texttt{proof} environment of 52 lines). No mathematics is written, repaired or re-derived: the statement and the proof are the article's own lines, in the article's own notation, and no retained line is substituted or re-flowed.  Retained uncounted as the source-local context the proof needs: lines 386--398.  Nothing was omitted from any retained span, so the package carries no omission record.  No retained line contains a citation, so the wrapper carries no bibliography and no bibliography entry.  No retained line contains a figure environment, an image-inclusion command or drawing code, so no image is delivered and the package declares no asset dependency.  Editorial formatting only, in addition: an AMS \texttt{amsart} preamble that restates the article's own declarations and macros used by the retained lines, namely the environments \texttt{lemma} on separate counters, together with the standard packages those lines need.  The retained environments print in this document as Lemma 1 in order of appearance, whereas the article prints the same items with the numbering of its own full document.  The complete native source of the article as distributed in the arXiv e-print for this exact version is bundled unchanged as \texttt{sources/199582/source0.tex}.
\begin{lemma}[One-cycle skew-product criterion]\label{lem:skew-product}
Let \(X\) be a finite set, \(A:X\to X\) a single cycle of length \(|X|\), and let
\[
F:X\times\mathbb Z_m\to X\times\mathbb Z_m,
\qquad F(x,y)=(A(x),y+\phi(x)).
\]
Put
\[
S=\sum_{x\in X}\phi(x)\in\mathbb Z_m.
\]
Then \(F\) is a single cycle on \(X\times\mathbb Z_m\) if and only if \(S\) generates
\(\mathbb Z_m\), equivalently \(\gcd(S,m)=1\).
\end{lemma}

\begin{lemma}[Prefix-count criterion]\label{lem:count-criterion}
Let \(W=(\xi_0,\ldots,\xi_{m-1})\) be a length-\(m\) word in \(\Sigma_6\).  For thresholds \(\tau_0,\ldots,\tau_{m-1}\), the associated return map is
\[
R_W=M_{\tau_{m-1}}^{\xi_{m-1}}\cdots M_{\tau_1}^{\xi_1}M_{\tau_0}^{\xi_0}.
\]
Let
\[
N_0,N_\Delta,N_2,N_3,N_4,N_5,N_6
\]
be the symbol counts of \(W\).  If
\[
\gcd(N_0,m)=1
\]
and
\[
\gcd(N_k-N_\Delta,m)=1\qquad(2\le k\le6),
\]
then the return map associated to \(W\) is a single \(m^6\)-cycle on \(Q_6\), for any choice of
thresholds in the layers.
\end{lemma}

\begin{proof}
We prove the more general statement in dimension \(r\le6\), allowing the additional full-translation
symbol \(T\).  Let the word length be \(m\).  Let \(N_T\) be the number of occurrences of \(T\).  The
criterion is:
\[
\gcd(m-N_0,m)=1,
\qquad
\gcd(N_k-N_\Delta,m)=1\quad(2\le k\le r).
\]
Since \(m-N_0\equiv -N_0\pmod m\), this is the same first condition as in the statement.  The symbol
\(T\) is counted among the nonzero symbols in \(m-N_0\), but it does not appear in the later
difference conditions.

For \(r=1\), all nonzero symbols, including the auxiliary symbol \(T\) when it occurs, act as the translation \(x\mapsto x-1\) on \(\mathbb Z_m\).  Thus the
return map is translation by minus the number of nonzero symbols, and it is a single cycle exactly when
\(m-N_0\) is a unit modulo \(m\).

Assume the statement known in dimension \(r-1\).  Write
\[
Q_r=Q_{r-1}\times\mathbb Z_m,
\qquad z=(u,y).
\]
Project each layer map to the \(u\)-coordinate.  The symbols \(0,\Delta,2,\ldots,r-1\) project to the
corresponding symbols in dimension \(r-1\).  The top numeric symbol \(r\) projects to the full
translation \(T_{r-1}:u\mapsto u-(1,\ldots,1)\).  The full translation symbol \(T\), if present, also projects to
\(T\).  Therefore the projected word satisfies the induction hypotheses, because the first condition is
unchanged and the difference conditions for \(2\le k\le r-1\) are unchanged.  Hence the projected
return map on \(Q_{r-1}\) is one cycle.

It remains to compute the drift in the last coordinate over one base cycle.  We apply Lemma~\ref{lem:skew-product} with \(X=Q_{r-1}\), with \(A\) equal to the projected return map, and with \(\phi(u)\) equal to the accumulated increment of the last coordinate during one full length-\(m\) word starting from base state \(u\).  The total drift \(S=\sum_u\phi(u)\) is the sum of the following contributions.

Consider a fixed occurrence of \(\Delta\) in the word.  Immediately before this occurrence, the base state is obtained from the initial base state by a bijection of \(Q_{r-1}\).  Hence, as the initial base point ranges over one base cycle, the state seen at this occurrence ranges over all of \(Q_{r-1}\) exactly once.  The \(\Delta\)-map decreases the last coordinate exactly when that state has no coordinate equal to the threshold used in that layer.  This count is independent of the threshold: for any fixed \(\tau\in\mathbb Z_m\), the number of points of \(Q_{r-1}\) with no coordinate equal to \(\tau\) is
\[
(m-1)^{r-1}.
\]
Thus each occurrence of \(\Delta\) contributes \(-(m-1)^{r-1}\) to the total drift.

A fixed occurrence of the top numeric symbol \(r\) similarly sees every base state once.  It decreases the last coordinate exactly when at least one coordinate of the base state equals the current threshold.  Hence each occurrence of \(r\) contributes
\[
-\bigl(m^{r-1}-(m-1)^{r-1}\bigr).
\]
The lower numeric symbols \(2,\ldots,r-1\) do not change the last coordinate.  A full translation symbol \(T\) decreases the last coordinate for every \(u\), hence contributes \(-m^{r-1}\equiv0\pmod m\) to the drift.

Thus the total drift in the last coordinate is congruent modulo \(m\) to
\[
-N_\Delta(m-1)^{r-1}-N_r\bigl(m^{r-1}-(m-1)^{r-1}\bigr)
\equiv (N_r-N_\Delta)(m-1)^{r-1}.
\]
Because \((m-1)^{r-1}\) is a unit modulo \(m\), Lemma~\ref{lem:skew-product} says that the final
fiber is one cycle exactly when \(N_r-N_\Delta\) is a unit modulo \(m\).  This is the required last
condition.  The induction is complete.
\end{proof}

\end{document}
